\section{第一阶段深讲：CPU 对操作系统暴露的契约}

CPU 不是操作系统可以随意想象的执行器。它给软件暴露的是一组精确契约：指令如何改变寄存器，异常在什么边界发生，特权模式如何切换，页表根在哪里，系统调用从哪里进入，返回用户态前哪些状态必须恢复。理解这些契约，比背一张 CPU 结构图更重要。

\subsection{寄存器不是临时变量，而是架构状态}

软件能直接观察和依赖的 CPU 状态称为架构状态。通用寄存器、程序计数器、栈指针、状态寄存器、控制寄存器、浮点和向量寄存器，都属于这个范围。微架构内部的乱序窗口、重命名寄存器、分支预测表通常不可直接观察，但它们会影响性能和安全。

\begin{longtable}{p{0.2\linewidth}p{0.34\linewidth}p{0.34\linewidth}}
\toprule
状态 & 对普通程序的意义 & 对 OS 的意义 \\
\midrule
RIP & 下一条或当前指令位置 & 异常返回、信号、调试、上下文切换 \\
RSP & 当前栈顶 & 用户栈、内核栈、ABI 对齐、栈溢出 \\
通用寄存器 & 参数、返回值、临时计算 & syscall ABI、trap frame、调度保存恢复 \\
RFLAGS & 条件码、中断使能、模式位 & 中断开关、返回用户态、单步调试 \\
页表根寄存器 & 当前地址空间 & 进程切换、内核映射、TLB 语义 \\
向量寄存器 & SIMD/浮点计算 & 信号帧、延迟保存、上下文切换成本 \\
\bottomrule
\end{longtable}

一条用户指令被中断打断后，程序将来还要像没发生过中断一样继续运行。为此，内核不能只保存 RIP，还要保存足够恢复用户可见状态的寄存器。x86-64 ABI 和入口代码共同定义保存集合；凡是用户态能依赖的状态，内核要么保存恢复，要么在 ABI 中明确它会被破坏。

\subsection{指令执行不是一句“取指执行”}

最小模型可以写成取指、译码、执行、访存、写回。真实 CPU 可能乱序、预测和并行，但最终必须提交出与 ISA 一致的结果。OS 关心的是每条指令在异常和中断面前的边界。

\begin{lstlisting}
instruction retirement view:
  fetch bytes from instruction address
  decode according to ISA
  read architectural operands
  execute operation or memory access
  if fault, report precise exception
  otherwise commit architectural result
\end{lstlisting}

“提交”这个概念很关键。乱序 CPU 可以提前执行后面的指令，但不能让用户态看到不符合顺序语义的架构状态。若一条 load 触发 page fault，内核必须知道 faulting instruction 是哪一条，之前的指令已经生效，之后的指令不能对架构状态产生不可回滚影响。否则 page fault handler 无法修复状态并重试。

\subsection{精确异常是虚拟内存的地基}

精确异常的意思是，异常发生时 CPU 能给出一个像顺序执行一样的状态点：异常指令之前的指令已经完成，异常指令和之后的指令没有提交最终架构效果。缺页、除零、非法指令、断点、系统调用，都依赖这个性质。

\begin{lstlisting}
if instruction at RIP faults:
  saved_rip = RIP of faulting instruction
  saved_error = reason and access type
  architectural state corresponds to before this instruction commits
  kernel may fix the cause and return to saved_rip
\end{lstlisting}

没有精确异常，COW 就很难实现。写 COW 页的 store 必须在真正写入共享物理页之前停下来，让内核复制新页并更新 PTE，再回到同一条 store。若 store 已经部分写入旧页，父子隔离就被破坏。硬件异常语义直接决定 OS 能实现哪些内存抽象。

\subsection{fault、trap、interrupt 的差别}

中文里常把异常、中断、陷入混着说，但写内核时必须分清事件来源和返回语义。

\begin{longtable}{p{0.16\linewidth}p{0.28\linewidth}p{0.28\linewidth}p{0.18\linewidth}}
\toprule
类型 & 来源 & 典型例子 & 返回位置 \\
\midrule
fault & 当前指令无法完成 & page fault、保护错误 & 常返回同一条指令重试 \\
trap & 当前指令主动进入内核或调试 & syscall、breakpoint & 常返回下一条指令 \\
interrupt & 外部异步事件 & timer、device IRQ、IPI & 返回被打断位置 \\
abort & 严重硬件或不可恢复错误 & machine check、总线错误 & 可能无法安全返回 \\
\bottomrule
\end{longtable}

这个分类影响内核入口保存什么状态。fault 要保留足够信息让指令重试；syscall 要按 ABI 把返回值写回指定寄存器；interrupt 要避免破坏被打断上下文；严重硬件错误可能只能记录日志、隔离页或 panic。

\subsection{用户栈和内核栈为什么要分开}

用户栈由用户程序控制，可能被破坏、未映射、接近栈底，甚至故意指向恶意地址。内核进入特权态后不能继续信任用户栈，否则中断帧、返回地址和内核局部变量都会落在不可信内存里。因此系统调用、中断和异常入口通常切换到内核栈。

\begin{lstlisting}
on privilege transition:
  CPU or entry code chooses kernel stack for current task
  save user RIP, RSP, RFLAGS and selected registers
  run kernel handler on trusted stack
  restore or validate state before returning to user mode
\end{lstlisting}

每个线程需要自己的内核栈，因为线程可能在系统调用中阻塞。若所有线程共享一个内核栈，线程 A 阻塞在内核里后，线程 B 进入内核就会覆盖 A 的调用链。内核栈大小有限，所以内核代码不能随意放巨大局部数组，也不能无限递归。

\subsection{系统调用入口是一段极谨慎的汇编}

系统调用入口处，C 运行环境还不完整。入口代码要处理用户寄存器、切换栈、建立内核可用的调用约定、屏蔽或允许中断、检查系统调用号、进入 C 函数分发。返回路径同样危险：返回用户态前要处理信号、调度、审计、单步调试和可能的抢占。

\begin{lstlisting}
syscall entry outline:
  save user register frame
  switch to trusted kernel stack if needed
  sanitize user-provided state
  dispatch syscall table entry
  store return value in ABI register
  check pending signal or reschedule flag
  restore user state
  return to user mode
\end{lstlisting}

这里不能有“差不多”。若入口漏保存一个用户可见寄存器，程序会随机坏；若没有清理某些特权状态，用户态可能继承内核状态；若返回路径错误恢复 RFLAGS，用户程序可能以错误中断状态或错误单步状态运行。

\subsection{上下文切换不是只保存寄存器}

最浅的说法是“上下文切换保存现场”。更准确地说，它要切换当前 CPU 正在代表哪个任务执行。这个切换至少涉及调度实体、内核栈、寄存器保存区、地址空间、抢占计数、FPU 状态和 CPU 本地 current 指针。

\begin{lstlisting}
switch from prev to next:
  prev.thread.rsp = current RSP
  save callee-saved registers into prev thread context
  current = next
  load next address-space root if mm changes
  switch RSP to next.thread.rsp
  restore next registers
  return as if next's old schedule call just finished
\end{lstlisting}

注意最后一句：上下文切换不是“跳到 next 的 main 函数”。next 之前可能正在系统调用中睡眠，也可能刚被定时器抢占。恢复后，它从自己上次离开 CPU 的内核位置继续运行，最终回到用户态或继续内核逻辑。

\subsection{FPU 和 SIMD 为什么让切换变贵}

现代向量寄存器很大。若每次切换都保存所有向量状态，成本高；若延迟保存，又要处理任务第一次使用 FPU/SIMD 时的异常、跨任务泄漏和安全补丁。许多系统在性能和安全之间选择 eager 策略，即每次切换都明确保存恢复用户可见向量状态。

\begin{longtable}{p{0.18\linewidth}p{0.36\linewidth}p{0.34\linewidth}}
\toprule
策略 & 好处 & 风险或成本 \\
\midrule
lazy save & 不使用 FPU 的任务不付成本 & 状态管理复杂，历史上暴露过安全问题 \\
eager save & 语义清晰，隔离更直接 & 每次切换固定成本增加 \\
按需扩展状态 & 只为使用的扩展保存空间 & 需要精确追踪任务使用的 CPU 特性 \\
\bottomrule
\end{longtable}

这解释了为什么 CPU 新增向量扩展不只是编译器问题。OS 要知道如何启用、保存、恢复、暴露给用户、在信号帧中表示，并在不支持时给出正确非法指令异常。

\subsection{控制寄存器把硬件策略交给内核}

页表根、向量表基址、系统调用入口、定时器配置、中断屏蔽、cache 策略、性能计数器控制，这些状态通常只能在内核或更高特权层修改。它们让 OS 能建立策略，也让错误配置变得危险。

\begin{lstlisting}
privileged configuration examples:
  set page table root for current address space
  install exception vector table
  configure syscall entry point
  program local timer deadline
  enable or disable selected interrupts
  configure performance monitoring counters
\end{lstlisting}

普通用户程序不能直接改这些寄存器。若能改页表根，它可以访问任意内存；若能改异常向量，它可以劫持系统调用；若能关闭定时器中断，它可以阻止抢占。特权级的价值，就体现在这些寄存器的写权限被硬件强制限制。

\subsection{x86-64 名字与共同原则}

本册不再把多个 ISA 混在一起讲。x86-64 的名字就是本书后续源码和汇编的默认语言：RIP/RSP/RFLAGS 表示执行现场，CR3 表示当前页表根，IDT 表示中断和异常入口，\code{syscall} 表示用户进入内核，Ring 3/Ring 0 表示用户态和内核态。重要的是把这些名字背后的角色讲清：有用户态和特权态，有异常入口，有地址翻译，有中断，有返回用户态的受控路径，有保存恢复上下文的要求。

\begin{longtable}{p{0.24\linewidth}p{0.28\linewidth}p{0.36\linewidth}}
\toprule
问题 & x86-64 名字 & OS 角色 \\
\midrule
用户进入内核 & \code{syscall} & 受控系统调用入口 \\
特权层级 & Ring 3/Ring 0 & 用户态和内核态隔离 \\
页表根 & CR3 & 当前地址空间翻译根 \\
异常入口 & IDT & 中断、异常和 trap 的分发入口 \\
返回用户态 & \code{sysretq}/\code{iretq} & 恢复用户 RIP/RSP/RFLAGS 和 CPL \\
\bottomrule
\end{longtable}

学习 OS 时要把“名字”和“角色”分开。CR3 不是普通变量，它是地址空间切换的硬件根；IDT 不是函数表，它是异常和中断的受控入口；\code{syscall} 不是函数调用，它是跨特权级入口。

\subsection{CPU 契约的最低验收}

读完本节，应该能回答以下问题：为什么 page fault 可以返回同一条指令重试；为什么系统调用不能使用用户栈保存内核状态；为什么上下文切换要考虑地址空间和向量寄存器；为什么不同架构的系统调用指令不同但 OS 模型相似；为什么特权寄存器必须由硬件限制写入。

如果这些问题能回答清楚，进程、线程、系统调用和虚拟内存就不再是孤立概念，而是 CPU 契约自然推出的软件结构。
